% Part: incompleteness
% Chapter: incompleteness-provability
% Section: first-incompleteness-thm

\documentclass[../../../include/open-logic-section]{subfiles}

\begin{document}

\olfileid{inc}{inp}{1in}

\olsection{પ્રથમ અપૂર્ણતા પ્રમેય}

હવે પ્રથમ અપૂર્ણતા પ્રમેયની
ગોડેલની મૂળ સાબિતી
વર્ણવી શકીએ. $\Th{T}$
એ અંકગણિતની ભાષાનું
વિસ્તરણ કરતી ભાષામાંનો
કોઈપણ સંગણનીય રીતે
સ્વયંસિદ્ધિત સિદ્ધાંત હોય
અને $\Th{T}$માં
$\Th{Q}$નાં સ્વયંસિદ્ધિઓ
સમાવેલાં હોય. ખાસ કરીને
તેનો અર્થ એ છે કે
$\Th{T}$ સંગણનીય
વિધેયો અને સંબંધોનું
નિરૂપણ કરે છે.

અગાઉ દલીલ કરી છે કે
સૂત્રો અને સાબિતીઓને
સંખ્યાઓ વડે યોગ્ય રીતે
સંકેતિત કરીએ તો સંબંધ
$\Prf[\Th{T}](x,y)$
સંગણનીય છે. અહીં
$\Prf[\Th{T}](x,y)$
ત્યારે અને માત્ર ત્યારે જ
સાચો છે જ્યારે $x$ એ
$\Th{T}$માં ગોડેલ સંખ્યા
~$y$ ધરાવતા !!{formula}ના
!!a{derivation}ની ગોડેલ
સંખ્યા હોય. હકીકતમાં
ગોડેલે ધ્યાનમાં રાખેલા
ચોક્કસ સિદ્ધાંત માટે
પોતાના લેખની ૪૫ વિધેયો
અને સંબંધોની યાદીથી
તેમણે આ સંબંધ આદિમ
પુનરાવર્તી છે એવું બતાવ્યું.
૪૫મો સંબંધ, $x B y$,
તેમની $\Th{T}$ની ખાસ
પસંદગી માટેનો
$\Prf[\Th{T}](x,y)$
જ છે. યાદ રાખો કે
ગોડેલ પોતાના લેખમાં
``પુનરાવર્તી'' કહે છે
ત્યાં આજે આપણે
``આદિમ પુનરાવર્તી''
કહીશું.

$\Prf[\Th{T}](x,y)$ સંગણનીય
હોવાથી $\Th{T}$માં તેનું
નિરૂપણ થાય છે. તેને
નિરૂપતા સૂત્ર માટે
$\OPrf[\Th{T}](x,y)$
લખીશું. $\OProv[\Th{T}](y)$
એ સૂત્ર
$\lexists[x][\OPrf[\Th{T}](x,y)]$
હોય. આ ગોડેલની યાદીનો
૪૬મો સંબંધ,
$\fn{Bew}(y)$, વર્ણવે છે.
ગોડેલ નોંધે છે તેમ,
આ એકમાત્ર એવો સંબંધ
છે જે ``પુનરાવર્તી છે''
એવું કહી શકાતું નથી.
કદાચ તેમનો અર્થ આ હતો:
વ્યાખ્યા પરથી તે સંગણનીય
છે એવું સ્પષ્ટ થતું નથી;
અને પછીનો વિકાસ બતાવે
છે કે હકીકતમાં તે
સંગણનીય નથી.

$\Th{T}$ એ~$\Th{Q}$ને
સમાવતો કોઈ
!!{axiomatizable} સિદ્ધાંત
હોય. ત્યારે
$\Prf[\Th{T}](x, y)$
નિર્ણેય છે અને તેથી
કોઈ !!a{formula}
~$\OPrf[\Th{T}](x, y)$
વડે~$\Th{Q}$માં નિરૂપાય
છે. ઉપર વર્ણવેલું
$\OProv[\Th{T}](y)$
સૂત્ર લો. નિશ્ચિત-બિંદુ
લેમ્માથી એવું સૂત્ર
$!G_\Th{T}$ છે કે
$\Th{Q}$ (અને તેથી
$\Th{T}$ પણ) નીચેનું
!!{derive} કરે છે:
\begin{equation}
\ollabel{eqn:qpf}
!G_\Th{T} \liff \lnot \OProv[\Th{T}](\gn{!G_\Th{T}}).
\end{equation}
નોંધો કે $!G_\Th{T}$નો
મૂળ અર્થ ``$!G_\Th{T}$
એ~$\Th{T}$માં
!!{derivable} નથી'' છે.

\begin{lem}\ollabel{lem:cons-G-unprov}
જો $\Th{T}$ એ~$\Th{Q}$નું
સુસંગત,
!!{axiomatizable} વિસ્તરણ
હોય, તો
$\Th{T} \Proves/ !G_\Th{T}$.
\end{lem}

\begin{proof}
માનો કે $\Th{T}$ એ
$!G_\Th{T}$ને !!{derive}
કરે છે. તો ખરેખર
!!a{derivation} \emph{છે};
તેથી કોઈ સંખ્યા~$m$ માટે
સંબંધ $\Prf[\Th{T}](m,
\Gn{!G_\Th{T}})$ સાચો છે.
પણ તેથી $\Th{Q}$ વાક્ય
$\OPrf[\Th{T}](\num m, \gn{!G_\Th{T}})$ને
!!{derive} કરે છે.
આથી $\Th{Q}$ એ
$\lexists[x][\OPrf[\Th{T}](x,\gn{!G_\Th{T}})]$
પણ !!{derive} કરે છે,
જે વ્યાખ્યા મુજબ
$\OProv[\Th{T}](\gn{!G_\Th{T}})$
છે. \olref{eqn:qpf}
મુજબ $\Th{Q}$ એ
$\lnot !G_\Th{T}$ને
!!{derive} કરે છે;
અને $\Th{T}$ એ
$\Th{Q}$નું વિસ્તરણ
હોવાથી $\Th{T}$ પણ
એમ કરે છે. બતાવ્યું કે
જો $\Th{T}$ એ
$!G_\Th{T}$ને !!{derive}
કરે, તો $\lnot !G_\Th{T}$ને
પણ !!{derive} કરે છે;
તેથી તે અસુસંગત બને.
\end{proof}

\begin{defn}
\ollabel{thm:oconsis-q}
સિદ્ધાંત $\Th{T}$ને
\emph{$\omega$-સુસંગત}
ત્યારે કહીએ જ્યારે આ
શરત સંતોષાય: જો
$\lexists[x][!A(x)]$
કોઈપણ વાક્ય હોય અને
$\Th{T}$ એ
$\lnot
!A(\num 0)$,
$\lnot !A(\num 1)$,
$\lnot !A(\num 2)$,
\dots બધાંને
!!{derive} કરે,
તો $\Th{T}$ એ
$\lexists[x][!A(x)]$ને
સાબિત કરતું નથી.
\end{defn}

દરેક $\omega$-સુસંગત
સિદ્ધાંત સુસંગત પણ છે.
કારણ કે જો $\Th{T}$
અસુસંગત હોય, તો દરેક
~$!A$ માટે
$\Th{T} \Proves !A$.
ખાસ કરીને, જો
$\Th{T}$ અસુસંગત હોય,
તો દરેક~$n$ માટે
$\lnot !A(\num n)$ને
!!{derive} કરે છે અને
$\lexists[x][!A(x)]$ને
પણ !!{derive} કરે છે.
તેથી $\Th{T}$ અસુસંગત
હોય તો તે
$\omega$-અસુસંગત છે.
વિપરિત અનુમાનથી,
$\Th{T}$ જો
$\omega$-સુસંગત હોય
તો તે સુસંગત હોવું જ
જોઈએ.

\begin{lem}\ollabel{lem:omega-cons-G-unref}
જો $\Th{T}$ એ~$\Th{Q}$નું
$\omega$-સુસંગત,
!!{axiomatizable} વિસ્તરણ
હોય, તો
$\Th{T} \Proves/ \lnot !G_\Th{T}$.
\end{lem}

\begin{proof}
જો $\Th{T}$ એ
$\lnot !G_\Th{T}$ને
!!{derive} કરે તો તે
$\omega$-અસુસંગત છે
એવું બતાવીએ. માનો કે
$\Th{T}$ એ
$\lnot !G_\Th{T}$ને
!!{derive} કરે છે.
જો $\Th{T}$ અસુસંગત
હોય તો તે
$\omega$-અસુસંગત છે,
એટલે કામ પૂરું.
નહીં તો $\Th{T}$
સુસંગત છે, તેથી
\olref{lem:cons-G-unprov}
મુજબ તે $!G_\Th{T}$ને
!!{derive} કરતું નથી.
$\Th{T}$માં
$!G_\Th{T}$નું
કોઈ !!{derivation} નથી;
તેથી $\Th{Q}$ નીચેનું
!!{derive} કરે છે:
\[
\lnot \OPrf[\Th{T}](\num 0, \gn{!G_\Th{T}}), \lnot \OPrf[\Th{T}](\num 1,
\gn{!G_\Th{T}}), \lnot \OPrf[\Th{T}](\num 2, \gn{!G_\Th{T}}), \dots
\]
અને $\Th{T}$ પણ એમ
કરે છે. બીજી બાજુ,
\olref{eqn:qpf} મુજબ
$\lnot
!G_\Th{T}$ એ
$\lexists[x][\OPrf[\Th{T}](x,\gn{!G_\Th{T}})]$
ની સમકક્ષ છે. તેથી
$\Th{T}$
$\omega$-અસુસંગત છે.
\end{proof}

\begin{prob}
  દરેક $\omega$-સુસંગત
  સિદ્ધાંત સુસંગત હોય છે.
  વિપરીત વિધાન ખોટું છે,
  એટલે સુસંગત છતાં
  $\omega$-અસુસંગત
  સિદ્ધાંતો છે, તે બતાવો.
  આ માટે $\Th{Q} \cup
  \{\lnot !G_\Th{Q}\}$
  સુસંગત પરંતુ
  $\omega$-અસુસંગત છે
  તે સાબિત કરો.
\end{prob}

\begin{thm}
\ollabel{thm:first-incompleteness}
$\Th{T}$ એ~$\Th{Q}$નું
કોઈપણ $\omega$-સુસંગત,
!!{axiomatizable} વિસ્તરણ
હોય. તો $\Th{T}$ પૂર્ણ
નથી.
\end{thm}

\begin{proof}
  જો $\Th{T}$
  $\omega$-સુસંગત છે,
  તો તે સુસંગત છે; તેથી
  \olref{lem:cons-G-unprov}
  મુજબ $\Th{T}
  \Proves/ !G_\Th{T}$.
  \olref{lem:omega-cons-G-unref}
  મુજબ
  $\Th{T} \Proves/ \lnot !G_\Th{T}$.
  આથી $\Th{T}$ અપૂર્ણ
  છે: તે $!G_\Th{T}$
  કે $\lnot !G_\Th{T}$,
  બેમાંથી એકેયને
  !!{derive} કરતું નથી.
\end{proof}

\end{document}
